\section{Generation by Denoising and Iterative Refinement}

The previous sections presented three major paradigms for generative modeling:
latent-variable decoding, adversarial generation, and exact transport by invertible maps.
Diffusion introduces a fourth viewpoint.
Instead of generating a sample in one shot, or evaluating density by a single formula, it models generation as the gradual reversal of corruption.

\medskip

\noindent
\textbf{Guiding idea.}
A complicated global distribution may be easier to learn through many small local denoising steps than through one direct global map.

\medskip

This section is intentionally slower and more conceptual than the next one.
Its goal is to make precise what is being corrupted, what is being learned, and why repeated denoising can become a generative mechanism.

\subsection*{From clean data to noisy data}

Suppose $x_0\in \mathbb{R}^d$ denotes a data point drawn from the unknown data distribution.
A denoising-based generative method begins by defining a family of corrupted versions
\[
x_1,x_2,\dots,x_T,
\]
of $x_0$, where larger $t$ corresponds to more corruption.
The variable $x_T$ should be so heavily corrupted that it is close to a simple reference distribution such as an isotropic Gaussian.

\begin{definition}[Forward corruption process]
A forward corruption process is a conditional distribution
\[
q(x_{1:T}\mid x_0)
\]
that maps clean data $x_0$ into progressively noisier variables $x_1,\dots,x_T$.
In diffusion models this process is usually chosen by design rather than learned from data.
\end{definition}

The choice to fix the forward process is important.
It means that the difficult modeling problem is shifted away from ``how should we corrupt data?'' and toward ``how should we reverse that corruption?''

\subsection*{A clean probabilistic formulation of gradual corruption}

A convenient formulation is a Markov chain,
\[
q(x_{1:T}\mid x_0)=\prod_{t=1}^T q(x_t\mid x_{t-1}),
\]
where each step adds only a small amount of noise.
A standard Gaussian choice is
\[
q(x_t\mid x_{t-1})=
\mathcal{N}\!\left(x_t\mid \sqrt{1-\beta_t}\,x_{t-1},\, \beta_t I\right),
\]
with $0<\beta_t<1$ typically small.

\medskip

This formula has a direct interpretation:
\begin{itemize}
    \item the factor $\sqrt{1-\beta_t}$ slightly shrinks the previous signal,
    \item the term $\beta_t I$ injects fresh Gaussian noise,
    \item repeated application gradually destroys fine structure.
\end{itemize}

For a single corruption level, one may also study a simpler conditional model of the form
\[
q_\sigma(\tilde x\mid x)=\mathcal{N}(\tilde x\mid x,\sigma^2 I),
\]
where $\tilde x=x+\sigma\varepsilon$ and $\varepsilon\sim \mathcal{N}(0,I)$.
This ``single-step corruption'' viewpoint is historically useful because it connects diffusion to denoising autoencoders.

\subsection*{Denoising as learning reverse conditional structure}

If corruption is probabilistic, then denoising is naturally a conditional inference problem.
Given a noisy observation $x_t$, one would like to understand the conditional distribution of the less noisy variable that preceded it.
This suggests learning reverse conditionals of the form
\[
p_\theta(x_{t-1}\mid x_t).
\]

\begin{definition}[Reverse denoising model]
A reverse denoising model is a parameterized family of conditional distributions
\[
p_\theta(x_{t-1}\mid x_t),\qquad t=1,\dots,T,
\]
intended to approximate the true reverse conditionals induced by the forward process.
\end{definition}

If these reverse conditionals were exact, then one could generate by starting from a simple sample $x_T$ and successively drawing
\[
x_{T-1}\sim p_\theta(x_{T-1}\mid x_T),\quad
x_{T-2}\sim p_\theta(x_{T-2}\mid x_{T-1}),\quad \dots,\quad
x_0\sim p_\theta(x_0\mid x_1).
\]
Thus generation becomes iterative refinement:
noise is converted into structure one small step at a time.

\paragraph{What is being learned?}
The reverse model does not need to memorize entire clean samples.
Rather, at each stage it must capture the local conditional geometry of the data distribution:
which slightly cleaner configurations are plausible given the current noisy one.

\subsection*{From denoising autoencoders to iterative generation}

Before diffusion models were formulated as full latent Markov chains, denoising autoencoders already used corruption as a learning device.
A denoising autoencoder samples a corrupted input
\[
\tilde x\sim q_\sigma(\tilde x\mid x),
\]
and trains a network $r_\theta(\tilde x)$ to reconstruct the clean target $x$ by minimizing a loss such as
\[
\mathcal{L}_{\mathrm{DAE}}(\theta)
=
\mathbb{E}_{x\sim p_{\mathrm{data}},\,\tilde x\sim q_\sigma(\tilde x\mid x)}
\bigl[\|x-r_\theta(\tilde x)\|^2\bigr].
\]

This is already a meaningful representation-learning problem:
the model must learn which directions in input space correspond to signal and which correspond to nuisance perturbation.
However, by itself, a denoising autoencoder is not yet a complete generative model in the same direct sense as a latent-variable sampler or a flow.
It learns local restoration, but not necessarily a full multi-step sampling scheme.

Diffusion extends this idea by:
\begin{itemize}
    \item defining a whole family of corruption levels rather than one fixed noise level,
    \item learning denoising rules indexed by noise scale or time step,
    \item chaining these rules into a generative process from nearly pure noise back to data.
\end{itemize}

\medskip

This is the conceptual bridge:
\begin{quote}
A denoising autoencoder learns to undo corruption. A diffusion model organizes many such denoising tasks into a stochastic path from noise to data.
\end{quote}

\subsection*{A one-step Gaussian denoising calculation}

The basic denoising idea can be seen in the simplest Gaussian setting.
Assume the clean variable is one-dimensional with prior
\[
X\sim \mathcal{N}(0,1),
\]
and the observed noisy variable is
\[
Y=X+\varepsilon,
\qquad \varepsilon\sim \mathcal{N}(0,\sigma^2),
\]
with $\varepsilon$ independent of $X$.
Then the conditional distribution $p(x\mid y)$ is also Gaussian.
Its mean is the Bayes-optimal denoiser under squared loss.

\begin{proposition}[Posterior mean for Gaussian denoising]
In the model above,
\[
X\mid Y=y \sim \mathcal{N}\!\left(\frac{1}{1+\sigma^2}y,\,\frac{\sigma^2}{1+\sigma^2}\right).
\]
In particular, the minimum mean-squared error estimator is
\[
\mathbb{E}[X\mid Y=y]=\frac{1}{1+\sigma^2}y.
\]
\end{proposition}

\begin{proof}[Proof sketch]
Because $(X,Y)$ is jointly Gaussian, the conditional distribution is Gaussian.
Using
\[
\operatorname{Var}(X)=1,
\qquad
\operatorname{Var}(Y)=1+\sigma^2,
\qquad
\operatorname{Cov}(X,Y)=1,
\]
the standard conditional Gaussian formula gives
\[
\mathbb{E}[X\mid Y=y]
=
\frac{\operatorname{Cov}(X,Y)}{\operatorname{Var}(Y)}y
=
\frac{1}{1+\sigma^2}y,
\]
and
\[
\operatorname{Var}(X\mid Y)
=
1-\frac{\operatorname{Cov}(X,Y)^2}{\operatorname{Var}(Y)}
=
\frac{\sigma^2}{1+\sigma^2}.
\]
\end{proof}

\paragraph{Interpretation.}
The optimal denoiser does not simply output $y$.
It shrinks $y$ toward the mean $0$, because the model knows that part of the observation is noise.
As $\sigma^2\to 0$, the estimator approaches $y$.
As $\sigma^2\to \infty$, the observation becomes uninformative and the estimator shrinks strongly toward $0$.

This toy example is small, but it captures the essential probabilistic logic of denoising:
learning a reverse conditional distribution means learning how clean structure is statistically related to noisy structure.

\subsection*{A worked iterative toy example}

Consider now a three-stage refinement picture in one dimension.
Suppose $x_0\sim \mathcal{N}(0,1)$, and define two corrupted versions
\[
x_1 = x_0 + \varepsilon_1,
\qquad
x_2 = x_1 + \varepsilon_2,
\]
with independent noises $\varepsilon_1\sim \mathcal{N}(0,0.25)$ and $\varepsilon_2\sim \mathcal{N}(0,1)$.
Then
\[
x_2 = x_0 + (\varepsilon_1+\varepsilon_2),
\]
so the total noise variance from $x_0$ to $x_2$ is $1.25$.
The optimal one-step estimator of $x_0$ from $x_2$ is therefore
\[
\mathbb{E}[x_0\mid x_2]
=
\frac{1}{1+1.25}x_2
=\frac{4}{9}x_2.
\]
Likewise, since $x_2=x_1+\varepsilon_2$ and $x_1\sim \mathcal{N}(0,1.25)$,
\[
\mathbb{E}[x_1\mid x_2]
=
\frac{1.25}{1.25+1}x_2
=\frac{5}{9}x_2.
\]
Then from $x_1$, the optimal denoiser back to $x_0$ is
\[
\mathbb{E}[x_0\mid x_1]
=
\frac{1}{1+0.25}x_1
=0.8x_1.
\]

If, for instance, we observe $x_2=1.8$, then the best intermediate denoised estimate is
\[
\hat x_1=\frac{5}{9}\cdot 1.8 = 1.0,
\]
and the next denoised estimate is
\[
\hat x_0=0.8\cdot 1.0=0.8.
\]
Direct denoising gives
\[
\mathbb{E}[x_0\mid x_2]=\frac{4}{9}\cdot 1.8 = 0.8
\]
as expected.

\paragraph{Why keep the multi-step view?}
In this simple Gaussian example, direct denoising and iterative denoising agree because everything is linear and exactly tractable.
But in complex high-dimensional data, the multi-step decomposition is useful because each local denoising task may be much simpler than the full jump from noise to data.
This is the practical intuition behind diffusion.

\subsection*{Noise prediction versus clean-sample prediction}

There are several equivalent-looking ways to parameterize a denoiser.
One may try to predict:
\begin{itemize}
    \item the clean sample $x_0$ from a noisy input,
    \item the previous state $x_{t-1}$ from $x_t$,
    \item or the actual noise that was added.
\end{itemize}

These choices are related, but they are not merely cosmetic.
Different parameterizations lead to different optimization behavior and different interpretations of what the network is learning.
In modern diffusion practice, predicting the added noise is particularly common because it interacts conveniently with the Gaussian forward process and leads to simple training objectives.
The next section makes this algebra explicit.

\subsection*{Progressive corruption and recovery}

\begin{figure}[h]
\centering
\fbox{%
\begin{minipage}{0.9\textwidth}
\centering
\vspace{0.2em}
\textbf{Progressive corruption and iterative recovery}\par\medskip
\begin{tabular}{c c c c c}
clean data & $\longrightarrow$ & lightly noisy & $\longrightarrow$ & heavily noisy \\
$x_0$ &  & $x_1$ &  & $x_T \approx \mathcal{N}(0,I)$
\end{tabular}

\medskip

\begin{tabular}{c c c c c}
$\mathcal{N}(0,I)$ & $\longrightarrow$ & coarse structure & $\longrightarrow$ & sample \\
$x_T$ &  & $x_{T-1},x_{T-2},\dots$ &  & $x_0$
\end{tabular}

\medskip

Forward process: fixed stochastic corruption.\\
Reverse process: learned conditional denoising steps.
\vspace{0.2em}
\end{minipage}}
\caption{Diffusion-style generation as reversal of a progressive corruption process.}
\end{figure}

The forward and reverse arrows should not be confused.
The forward process is usually chosen analytically and known exactly.
The reverse process is what must be learned from data.

\subsection*{What is already understood, and what remains heuristic}

It is useful to separate three levels of claim.

\paragraph{Probabilistically well defined.}
Once a forward corruption process is specified, the induced reverse conditionals are mathematically meaningful objects.
The question ``what is the true distribution of a slightly cleaner sample given a noisier one?'' is a legitimate conditional inference question.

\paragraph{Conceptually persuasive.}
The idea that complex generation can be decomposed into many local denoising problems is strongly supported by modern practice.
It explains why diffusion can be more stable than adversarial training and why denoising losses can be easier to optimize.

\paragraph{Still design practice.}
Specific noise schedules, parameterizations, samplers, and acceleration tricks often come from empirical design rather than from a single complete theory that tells us the unique correct choice.
Thus diffusion is mathematically structured, but not fully determined by theory alone.

\subsection*{What to remember}

\begin{itemize}
    \item Diffusion starts from a forward corruption process and learns to reverse it.
    \item Denoising is a conditional inference problem, not merely a heuristic cleanup step.
    \item Iterative refinement turns many local denoising tasks into a generative procedure.
    \item Denoising autoencoders provide the conceptual bridge to full diffusion models.
\end{itemize}

\subsection*{Common confusion}

\begin{itemize}
    \item The forward process is usually fixed, not learned.
    \item A denoiser is not automatically a full generative model unless it is embedded into a sampling procedure.
    \item Predicting noise, predicting $x_0$, and predicting $x_{t-1}$ are related but not identical parameterizations.
\end{itemize}

\subsection*{Exercises}

\begin{enumerate}
    \item Let $X\sim \mathcal{N}(0,1)$ and $Y=X+\varepsilon$ with $\varepsilon\sim \mathcal{N}(0,4)$. Compute $\mathbb{E}[X\mid Y=y]$ and $\operatorname{Var}(X\mid Y)$. Interpret the result.
    \item Suppose $q_\sigma(\tilde x\mid x)=\mathcal{N}(\tilde x\mid x,\sigma^2 I)$. Explain why minimizing $\mathbb{E}\|x-r_\theta(\tilde x)\|^2$ trains the posterior mean estimator when the model class is rich enough.
    \item Compare one-step denoising and multi-step denoising conceptually. Why might the multi-step version be easier in high dimensions?
    \item In your own words, explain the conceptual difference between a denoising autoencoder and a diffusion sampler.
\end{enumerate}

\subsection*{Notes and References}

Denoising as a learning principle predates modern diffusion models and appears prominently in work on denoising autoencoders and score-based viewpoints.
Modern diffusion models combine this denoising intuition with a carefully designed forward noising process and a learned reverse chain.
Historically important reference points include denoising autoencoders, score matching ideas, and the later diffusion and score-based generative-model literature.
The next section develops the formal diffusion equations that make these ideas computationally effective.
